% ECONOMICS_PROBLEM: {"id": "L41-549", "title": "Merger prediction accuracy", "jel_code": "L41", "source_id": "OP-0392"}
% !TeX program = xelatex
\documentclass[11pt,a4paper]{article}
\usepackage[margin=23mm]{geometry}
\usepackage{fontspec}
\setmainfont{Latin Modern Roman}
\usepackage{amsmath,amssymb,mathtools}
\usepackage{xcolor,enumitem,booktabs,longtable,array,xurl,hyperref}
\definecolor{muted}{HTML}{53636C}
\hypersetup{colorlinks=true,urlcolor=blue,linkcolor=blue}
\setlength{\parindent}{0pt}
\setlength{\parskip}{5pt}
\newcommand{\R}{\mathbb R}
\newcommand{\E}{\mathbb E}
\newcommand{\N}{\mathbb N}
\newcommand{\ind}{\mathbf 1}
\newcommand{\Var}{\operatorname{Var}}
\newcommand{\Cov}{\operatorname{Cov}}
\newcommand{\supp}{\operatorname{supp}}
\DeclareMathOperator*{\argmin}{arg\,min}
\DeclareMathOperator*{\argmax}{arg\,max}
\newcommand{\entrymeta}[3]{{\sffamily\small\color{muted}\textbf{\texttt{#1}}\quad|\quad #2\par #3\par}}
\newcommand{\Problemref}[1]{\texttt{#1}}
\newcommand{\JELref}[2]{\texttt{#2} (JEL \texttt{#1})}
\begin{document}
\section*{Merger prediction accuracy}
\textbf{Problem ID:} \texttt{L41-549}\quad \textbf{JEL:} \texttt{L41}\par
% BEGIN ECONOMICS STATEMENT
\label{problem:OP-0392}
{\sffamily\small\color{muted}Original subject group: Industrial organization and competition\par}
\label{definitions:problem:30007674}
\textbf{Audit classification:} Published research agenda. Blocked or abandoned mergers cannot be scored against an observed merger treatment effect without an additional design. Price denominators, weights and product discontinuation required clarification.
\subsection*{Definitions and precise research target}
Consider horizontal mergers in United States consumer-packaged-goods markets that close within a prespecified five-year cohort. Index mergers by \(m=1,\ldots,M\), with \(M\) the cohort size. Let \(X_m\) contain premerger shares, estimated substitution patterns, ownership, costs and observable entry conditions. A screening model \(f_k\), where \(k\) indexes a prespecified finite set of models, predicts the percentage price effect at horizon \(h=24\) months using only \(X_m\). Define the causal effect \(\tau_m(h)\) as the expenditure-weighted percentage difference between equilibrium prices with the merger and prices without it; both regimes permit rivals to reprice and entrants to respond. Cohort expenditure weights are nonnegative and sum to one. The prediction target is \(R_k=E[(f_k(X_m)-\tau_m(h))^2]\), where expectation averages over the cohort's merger distribution. Develop a retrospective design identifying \(\tau_m(h)\) from comparable unaffected markets or justified structural restrictions, and separately propagate uncertainty in these retrospective estimates into \(R_k\). Evaluate models on mergers excluded from model construction, and use blocked or abandoned deals to inform no-merger counterfactuals only under explicit selection and extrapolation assumptions. Determine which observable circumstances predict systematic forecast error. Accuracy measured against raw before-after prices is insufficient because unrelated demand and cost changes also move prices. The target is predictive reliability under specified equilibrium adjustment, not a universal claim that any one screening model succeeds everywhere.

\textbf{Definitions, assumptions and mathematical qualifications.} Fix positive counterfactual baseline prices and a common basket of premerger product equivalents, with weights \(w_{jm}\geq0\), \(\sum_jw_{jm}=1\). Define \(\tau_m=100\sum_jw_{jm}\{p_{jm}(1,h)/p_{jm}(0,h)-1\}\). Discontinued products require a prespecified equivalent-quality replacement price or a separate extensive-margin outcome. A finite cohort has \(R_k=\sum_m\nu_m(f_k(X_m)-\tau_m)^2\), where \(\nu_m\geq0\), \(\sum_m\nu_m=1\); a sampling distribution replaces this sum when generalizing beyond the cohort. A causal retrospective design specifies its comparison group, treatment timing, spillover restrictions and conditional untreated trend equality. Forecasts are fixed using training data independent of the evaluation cohort. A blocked deal reveals a no-merger outcome, not its unobserved merger price effect.
\subsection*{Variants and assumption changes}
\begin{enumerate}
\item \textbf{Editorial, independent noisy benchmark.} If \(\widehat\tau_m\) is conditionally unbiased with independently estimated variance \(v_m\), then use \((f_k(X_m)-\widehat\tau_m)^2-v_m\) for an unbiased risk contribution; state independence conditions.
\item \textbf{Editorial, selected deals.} Evaluate approved and blocked proposals jointly only through a model or design identifying both regimes. Bound risk over feasible \(\tau_m\) intervals when approval selection prevents point identification.
\end{enumerate}
\subsection*{References and source audit}
\textbf{Checked source.} 24 August 2024 author draft, Section 6, printed p. 37, sixth research direction. Full primary text retrieved. \url{https://web.stanford.edu/~ayurukog/trendsincompetition.pdf}.
\begin{enumerate}\raggedright
\item\hypertarget{definitions:source:30007674:1}{} Carl Shapiro; Ali Yurukoglu. 2024. \emph{Trends in Competition in the United States: What Does the Evidence Show?}. 24 August 2024 author draft, Section 6, printed p. 37, sixth research direction.
\url{https://web.stanford.edu/~ayurukog/trendsincompetition.pdf}.
DOI: \url{https://doi.org/10.3386/w32762}.
\end{enumerate}

\subsection*{Common conventions: Source mathematical conventions}

These conventions apply to each entry unless explicitly replaced.

\textbf{Models and identification.} A primitive model \(m\) specifies all probability laws, feasible choices, information, equilibrium and measurement rules used by its target. The admissible class \(\mathcal M\) is fixed independently of outcomes. If \(P_O(m)\) is its observable probability law and \(\tau(m)\) the target, its identified set at observable law \(P\) is
\[
 \mathcal I_\tau(P)=\{\tau(m):m\in\mathcal M,\ P_O(m)=P\}.
\]
Point identification means that this set is a singleton. An empty set signals incompatibility of data and assumptions. Coordinate bounds need not describe the joint set. Bounds are sharp only if every retained target is attainable or arbitrarily approachable by compatible models. Finite bounds require explicit restrictions on unobserved quantities; unrestricted confounding, hidden wealth, extrapolation or arbitrary equilibrium selection can yield uninformative sets.

\textbf{Causal interventions.} A potential outcome \(Y_i(p)\) is the outcome under a complete policy regime \(p\), including timing, financing, information and market spillovers. The target population and units are fixed before treatment. With interference, use the complete assignment vector or a justified exposure mapping; individual no-interference notation alone cannot identify an economy-wide intervention. A difference of mean potential outcomes does not require the cross-world joint distribution. Conditioning on post-treatment migration, survival, enrollment or employment changes the estimand and needs separate justification.

\textbf{Statistical statements.} An estimator, confidence set, forecast or test specifies a sampling law, model class and loss/coverage criterion. Uniform \((1-\alpha)\)-coverage means \(\inf_{m\in\mathcal M}\Pr_m\{\tau(m)\in C_n\}\geq1-\alpha\), or an explicitly asymptotic version. The whole parameter space trivially has coverage, so informativeness requires a stated width, risk or power objective. A forecasting improvement is relative to a named benchmark, information set and evaluation population.

\textbf{Equilibrium and optimization.} An equilibrium comprises feasible plans, optimality, beliefs where needed, budget constraints and all market/resource clearing. The primitives or a stated benchmark define each correspondence \(\mathcal E(p;m)\). Nonemptiness is assumed or proved, never inferred just from compact policy sets. A selected-equilibrium target uses a declared selection rule; a robust target quantifies over all permitted equilibria. Use suprema or infima unless attainment is established. An \(\varepsilon\)-optimal policy is feasible and within \(\varepsilon\) of the finite optimum value. Report an empty feasible set separately.

\textbf{Welfare and accounting.} Normative utility, interpersonal normalization, social weights, numeraire, discounting and the use of public receipts are primitives. Behavioral utility need not coincide with the welfare criterion. Household welfare includes ownership income and rents once; adding profits again to compensating variation generally double-counts them. A monetary equivalent is defined by an explicit transfer and price convention, with monotonicity and range conditions. Average effects, mechanism contributions, transfers and welfare losses are different targets. Infinite sums require convergence and dynamic feasible plans require terminal or no-Ponzi conditions.

\textbf{Source versions.} A working-paper date, revision date and journal publication year are distinguished. A DOI for a journal version is not automatically the DOI of an earlier manuscript. Locators refer to the stated version and printed pages where specified. A metadata-only check does not verify a claim about a full-text conclusion. Every entry's source subsection narrows the provenance claim accordingly.
% END ECONOMICS STATEMENT
\end{document}
